\documentclass[11pt]{exam}

\usepackage{amsfonts, amsmath, amssymb, amsthm}
\usepackage{mathtools}
\usepackage{algpseudocode}
\usepackage{geometry}
\geometry{margin=1in}
\usepackage{algpseudocode}
\usepackage{lmodern}

\renewcommand{\alglinenumber}[1]{#1}
\algtext*{EndFor}
\algtext*{EndWhile}
\algtext*{EndIf}

%% ===> Macros <===
\newcommand\N{\mathbb{N}}
\newcommand\Z{\mathbb{Z}}
\newcommand\Q{\mathbb{Q}}
\newcommand\R{\mathbb{R}}
\newcommand\C{\mathbb{C}}
\newcommand\E{\mathbb{E}}
%% ===> END Macros <===

%% ===> Deliminaters <===
\DeclarePairedDelimiter\ceil{\lceil}{\rceil}
\DeclarePairedDelimiter\floor{\lfloor}{\rfloor}
\DeclarePairedDelimiter\anglebrac{\langle}{\rangle}
\DeclarePairedDelimiter\norm{\lVert}{\rVert}
%% ===> END Deliminaters <===

\newcommand{\pset}[1]{\mathcal{P}(#1)}
\newcommand{\nats}{\mathbb{N}}
\newcommand{\natsp}{\mathbb{N}^+}
\newcommand{\ints}{\mathbb{Z}}
\newcommand{\intsp}{\mathbb{Z}^+}
\newcommand{\rats}{\mathbb{Q}}
\newcommand{\reals}{\mathbb{R}}
\newcommand{\realsp}{\mathbb{R}^+}
\newcommand{\cross}{\times}
\newcommand{\bins}{\{0,1\}}

\begin{document}

\title{CSC373 Assignment 2}
\author{
  Seokjin Yoo (1011093264) \\
  \texttt{hyde.yoo@mail.utoronto.ca}
  \and
  Troy Woodley (1011709705) \\
  \texttt{troy.woodley@mail.utoronto.ca}
}
\maketitle

\section{Question 1}
\subsection{Part A}

We define a procedure called \textsc{SomeFlow}, which takes in the
flow network $G=(V,E)$ and the lower bounds $\ell(e)$ for every edge
$e \in E$. The procedure finds two paths, one from $s$ to $u$ and
another from $v$ to $t$, for each edge $(u,v) \in E$, and increments
the flow of all edges on both paths and $(u,v)$ by $1$.

\vspace{\baselineskip}
\noindent\textsc{SomeFlow}$( G , \ell )$:
\begin{algorithmic}[1]
  \State Initialize $ f(e) = 0 $ for all $e \in E$
  \For{each edge $ (u, v) $ in $ E $} \label{1aouter1}
  \State Take $ p_1, p_2 $ to be a path from $ s $ to $ u $ and from
  $ v $ to $ t $ \label{1ainner1} \Comment{Using BFS}
  % from $ s $ and from $ t $ on $ G $ with all of its edges reversed
  \For{each edge $ e $ in $ p_1 $} \label{1ainner2}
  \State $ f(e) = f(e) + 1 $
  \EndFor
  \For{each edge $ e $ in $ p_2 $} \label{1ainner3}
  \State $ f(e) = f(e) + 1 $
  \EndFor
  \State $ f(u, v) = f(u, v) + 1 $ \label{1ainner4}
  \EndFor
  \State $ c = \max{\ell} $ \label{1aouter2}
  \For{each edge $ e $ in $ E $} \label{1aouter3}
  \State $ f(e) = f(e) * c $
  \EndFor
  \State\Return $ f $
\end{algorithmic}
\vspace{\baselineskip}

\noindent \textbf{Correctness}\\

First, consider the trivial case where all lower bounds $ \ell(e) = 0
$. If so, $c = 0$ from line 9, hence the final flow returned by
\textsc{SomeFlow} is the trivial flow where all $f(e) = 0$ for all $e
\in E$. This satisfies flow conservation and the lower bound, hence
is a valid flow.

To prove correctness of the algorithm, we want to show inductively
that the flow $f$ from the algorithm satisfies flow conservation
after the $i$th iteration of the outer loop at line \ref{1aouter1}.

When $i = 0$ or before the first iteration, $f(e) = 0$ for all $e
\in E$ hence flow conservation holds trivially. More precisely, for
any vertex $v \in V \setminus \{s, t\}$, we have $f^\mathit{in} (u) =
f^\mathit{out} (u) = 0$.

Now suppose that the flow $f$ satisfies flow conservation after
iteration $i-1$. From line \ref{1ainner1}, the paths $p_1$ and $p_2$
exist by the problem assumption.
The algorithm then increments the flow for all edges on $p_1$, $p_2$,
and the edge $(u,v)$ (lines \ref{1ainner2} to \ref{1ainner4}).
Consider some vertex $w_1 \in (p_1 \setminus \{s,u\})$ and $w_2 \in (p_2
\setminus \{v,t\})$. Then notice after the first inner loop at line
\ref{1ainner2}, both $f^\mathit{in}(w_1)$ and $f^\mathit{out}(w_1)$
are incremented by 1. Similarly, after the second inner loop at line
\ref{1ainner3}, both $f^\mathit{in}(w_2)$ and $f^\mathit{out}(w_2)$
are incremented by 1.

Also, notice that if $u \neq s$, the incoming flow at $u$ from $p_1$
is incremented by 1. Since $f(u,v)$ is also incremented, the outgoing
flow at $u$ is also incremented by 1. Similarly, if $v \neq t$, the
incoming flow at $v$ from $(u,v)$ and outgoing flow from $p_2$ are
both incremented by 1.

If $u$ or $v$ is also an intermediate vertex of the other path, then
both its flow in and flow out are incremented by another 1, as shown above.
For any other vertex not in $p_1$ or $p_2$ and is not $u$ or $v$, its
flow in and flow out are not changed.

By the induction hypothesis, $f^\mathit{in}(w) = f^\mathit{out}(w)$
for every $w \in V \setminus \{s, t\}$. Since both quantities
increase by the same amount, flow conservation
is preserved after iteration $i$.

Therefore, after the first outer loop, the algorithm correctly
constructs a flow satisfying flow conservation. Moreover, since the
loop considers every edge $(u, v)$ once and each iteration increments
the edge's flow by at least 1 from line \ref{1ainner4}, we must have
$f(u, v) \geq 1$ for all $(u, v) \in E$.

Now, from line \ref{1aouter2}, $c = max_{e \in E} \ell (e)$. Since
all $f(e) \geq 1$ by construction, $c * f(e) \geq c \geq \ell(e)$.
Therefore, all edges satisfy its lower bound after the second outer
loop at line \ref{1aouter3}. Finally, flow conservation still holds
because $f^\mathit{in} (v) = f^\mathit{out} (v)$ from above, meaning
$c * f^\mathit{in} (v) = c * f^\mathit{out} (v)$.

\noindent \textbf{Runtime}\\

The first outer loop on line \ref{1aouter1} performs $m = |E|$
iterations. For each iteration of the first outer loop, line
\ref{1ainner1} performs two BFS to find the path from $ s $ to $ u $
and from $ v $ to $ t $, which takes $ O(m + n) $. Then, the inner
loops at lines \ref{1ainner2} and \ref{1ainner3} performs at most $
O(m) $ iterations of a constant time operation each. Hence, the first
loop's runtime is $ O(m \cdot (m + n + m + m)) = O(m^2 + mn) $. Then,
the maximum at line \ref{1aouter2} takes $ O(m) $ time. Similarly,
the second outer loop at line \ref{1aouter3} performs $ O(m) $
iterations of a constant time operation, hence takes $ O(m) $ time.

Therefore, the runtime of \textsc{SomeFlow} is $ O(m^2 + mn + m + m)
= O(m^2 + m n) $, which is polynomial in $m$ and $n$.

\subsection{Part B}

First, construct a new ``residual'' network $G'$. For every edge
$(u,v) \in E$, the forward edge $(u,v)$ will have capacity
$\mathit{val}(f_0)$ and the backwards edge $(v,u)$ will have capacity
$f_0(u,v)-\ell(u,v)$. Then, perform Ford-Fulkerson on the network
$G'$ to find the maximum flow $f'$ from $t$ to $s$.

Finally, construct the flow $f$ by calculating
\[
  f(u,v) = f_0(u,v) + f'(u,v)
  - f'(v,u).
\]

\noindent \textbf{Correctness}\\

We want to show that the resulting flow $f$ is a valid flow
that satisfies the lower bound and flow conservation. Then, we will
show that $f$ indeed is a minimum valid flow.

We know from construction, that $f'(u, v) \geq 0$ and $f'(v, u) \leq
f_0 (u, v) - \ell (u, v) \Leftrightarrow f_0 (u, v) \geq \ell(u, v) +
f'(v, u)$. So, $f_0 (u, v) + f'(u, v) - f'(v, u) \geq \ell(u, v) +
f'(v, u) + 0 - f'(v, u) = \ell(u, v)$.
Thus $f (u, v) \geq \ell(u, v)$ meaning lower bounds are satsified.

Also, for each vertex $x \in V \setminus \{s, t\}$, we know
${f'}^\mathit{in} (x) = {f'}^\mathit{out} (x)$ by correctness of
Ford-Fulkerson. From part A, we also know that $f_0^\mathit{in} (x) =
f_0^\mathit{out} (x)$.
Since a forward edge in $G'$ increases the flow in $G$ and a backward
edge decreases the flow in $G$, we have
\[
  f^\mathit{out}(x)-f^\mathit{in}(x) =
  (f_0^\mathit{out}(x)-f_0^\mathit{in}(x)) +
  ({f'}^\mathit{out}(x)-{f'}^\mathit{in}(x)).
\]
Since both terms are zero by flow conservation of $f_0$ and $f'$,
flow conservation also holds for $f$. Note that this also tells us
that $\mathit{val}(f) = \mathit{val}(f_0) - \mathit{val}(f')$.

For minimality of $f$, suppose for contradiction, that there exists
another valid flow $g$ in $G$ such that $\mathit{val}(g) < \mathit{val}(f)$.
Now, construct another flow $h$ in $G'$ based on the difference
between $f_0$ and $g$. For every edge $(u,v) \in E$, if $g(u,v) <
f_0(u,v)$, assign $f_0(u,v)-g(u,v)$ to the backward edge $(v,u)$.
Otherwise, assign $g(u,v)-f_0(u,v)$ to the forward edge $(u,v)$.

For the backward edges, we know that $g(u,v) \geq \ell(u,v)$, so
$f_0(u,v)-g(u,v) \leq f_0(u,v)-\ell(u,v)$. Hence, the backward
capacity constraints are satisfied. Furthermore, since both $f_0$ and
$g$ satisfy flow conservation, their differences also satisfy flow
conservation. Thus, $h$ satisfies flow conservation and has value
\[
  \mathit{val}(h) = \mathit{val}(f_0)-\mathit{val}(g).
\]

Note that if the construction of $h$ leads to a cycle, we can remove cycles from
$h$ without changing its value. After removing such cycles, every edge
carries at most $\mathit{val}(h)$ flow. Since $\mathit{val}(g) \geq 0$, we have
\[
  \mathit{val}(h) \leq \mathit{val}(f_0).
\]
Therefore, the forward capacity constraints are also satisfied,
giving a valid flow $h$ in $G'$.

Now, notice $\mathit{val}(f_0) - \mathit{val}(g) > \mathit{val}(f_0) -
\mathit{val}(f)$. Since we know from above that $\mathit{val}(f_0) -
\mathit{val}(f) = \mathit{val}(f')$, this means that
$\mathit{val}(f_0) - \mathit{val}(g) > \mathit{val}(f')$.

This is a contradiction, since $f'$ is a maximum flow in $G'$ by
correctness of Ford-Fulkerson, but we have constructed another valid
flow $h$ in $G'$ such that $\mathit{val}(h) > \mathit{val}(f')$.

Therefore, $f$ must be a minimum flow, as desired.

\noindent \textbf{Runtime}\\

Construction of the new network $G'$ requires calculation of
capacities of forward and backward edges, which takes $O(m)$ time.
Then, Ford-Fulkerson with fewest edge (shortest path) heuristic takes
$O(m^2 n)$. Finally, the construction of $f$ takes $O(m)$, as
calculating the minimum flow for each edge takes $O(1)$ time.

Therefore, the runtime of the above algorithm is $ O(m^2 n) $, which
is polynomial in $m$ and $n$.

\subsection{Part C}

\section{Question 2}

\section{Question 3}

\subsection{Part A}

\subsection{Part B}

\subsection{Part C}
\end{document}
